% =============================================================================
%  lec04.tex —— 第 4 讲 拟范畴（底本 PDF 第 28–38 页）
% =============================================================================
\lecture{4}{拟范畴}
\begin{center}\small\color{black!55}（底本 PDF 第 28–38 页）\end{center}

就本课程的目的而言，本讲不给出证明，也不按逻辑顺序陈述结论。

\section{拟范畴}

\nitem\label{L4-1.1} 上次课我们构造了 Quillen 等价
\[
|{-}|\colon \Set_\Delta \rightleftarrows \Top \colon \Sing.
\]
由此可得：
\begin{itemize}[leftmargin=2.2em,itemsep=0.1em]
\item 任何 $\infty$-群胚——由同伦假设，它可以等同于一个同伦型——也可以实现为 $\Set_\Delta$ 中的双纤维对象，即 Kan 复形。
\item $\infty$-群胚之间的任何函子都可以实现为相应 Kan 复形之间态射的同伦类。这里，两个态射 $f,g\colon X\to Y$ 同伦，当且仅当存在态射 $H\colon X\times\Delta^1\to Y$，使得 $H(-,0)=f$ 且 $H(-,1)=g$。
\end{itemize}

\nitem\label{L4-1.2} 回顾：单纯集 $X$ 是 Kan 复形，当且仅当它对所有 $n\ge 2$ 的角嵌入
\[
\begin{tikzcd}[column sep=3.4em,row sep=2.6em]
\Lambda^n_i \arrow[r] \arrow[d] & X \\
\Delta^n \arrow[ur, dashed] &
\end{tikzcd}
\]
具有右提升性质。下面我们用高阶范畴论的语言，考察不同 $(n,i)$ 选择下的这些提升性质。

\nitem\label{L4-1.3} 当 $n=2$、$i=0$ 时，提升性质说的是：给定态射 $f\colon x_0\to x_1$，任何形如 $x_0\to x_2$ 的态射都可以通过 $f$ 分解。当 $n=3$、$i=0$ 时，提升性质说的是：给定态射 $f\colon x_0\to x_1$，任何交换图
\[
\begin{tikzcd}[column sep=3.2em,row sep=2.2em]
& x_2 \arrow[dr] & \\
x_0 \arrow[ur] \arrow[rr] & & x_3
\end{tikzcd}
\]
——或者更确切地说，见证这样一个图的可逆 $2$-态射——都可以通过 $f$ 分解。这意味着上述 $2$-态射等价于一个 $2$-态射
\[
\begin{tikzcd}[column sep=3.2em,row sep=2.2em]
& x_2 \arrow[dr] & \\
x_1 \arrow[ur] \arrow[rr] & & x_3
\end{tikzcd}
\]
与恒等 $2$-态射 $\id_f$ 的横复合。

\nitem\label{L4-1.4} 对偶地，$i=n$ 时的提升性质说的是：给定态射 $f\colon x\to y$，任何以 $y$ 为末顶点的同伦连贯图都可以通过 $f$ 分解。同样，这是 $\infty$-群胚而非一般 $\infty$-范畴的特征：它蕴含给定的态射 $f$ 有右逆。

\nitem\label{L4-1.5} 另一方面，$0<i<n$ 时的提升性质说的是：给定对象 $x$，可以把任何以 $x$ 为末顶点的同伦连贯图与另一个以 $x$ 为首顶点的同伦连贯图复合起来。任何像范畴的数学实体都应当满足这一要求。

\nitem\label{L4-1.6} 上述讨论促使我们作如下定义：

\begin{definition}\label{L4-1.7}
（Boardman--Vogt，1973）单纯集 $X$ 称为\textbf{拟范畴}（quasi-category），若它对内角嵌入
\[
\begin{tikzcd}[column sep=3.4em,row sep=2.6em]
\Lambda^n_i \arrow[r] \arrow[d] & X \\
\Delta^n \arrow[ur, dashed] &
\end{tikzcd}
\qquad (0<i<n)
\]
具有右提升性质。
\end{definition}

\yizhu{拟范畴又称弱 Kan 复形：这一概念由 Boardman 与 Vogt 在 1973 年引入，1980 年代由 Joyal 系统发展，并命名为拟范畴（quasi-category）；Joyal 于 2002 年证明了拟范畴正是 $(\infty,1)$-范畴的一类模型。\textbf{[来源：A：Boardman--Vogt, \emph{Homotopy Invariant Algebraic Structures on Topological Spaces}, 1973；Joyal, \emph{Quasi-categories and Kan complexes}, 2002]}}

设 $X$ 是拟范畴。我们把 $0$-单形 $x\in X_0$ 称为 $X$ 的\textbf{对象}，把 $1$-单形 $e\in X_1$ 称为从 $d^1_0(e)$ 到 $d^1_1(e)$ 的\textbf{态射}。

设 $X$ 与 $Y$ 是拟范畴，我们把 $\Set_\Delta$ 中的态射 $F\colon X\to Y$ 称为从 $X$ 到 $Y$ 的\textbf{函子}。

设 $\mathrm{QCat}\subset\Set_\Delta$ 为拟范畴构成的全子范畴。

\begin{example}\label{L4-1.8}
Kan 复形是拟范畴。
\end{example}

\begin{exercise}\label{L4-1.9}
范畴的脉是拟范畴。
\end{exercise}

\begin{exercise}\label{L4-1.10}
拟范畴的反范畴 $X^{\op}$ 是拟范畴。（见第 3 讲构造 1.14。）
\end{exercise}

\nitem\label{L4-1.11} 1980 年代，Joyal 提出：
\begin{center}
\emph{拟范畴是 $(\infty,1)$-范畴的模型。}
\end{center}
本讲把上述提法当作既定事实，并用它来发展关于 $(\infty,1)$-范畴的基本语言。我们将在下一讲为这一提法提供依据。

\section{态射的复合}

\begin{definition}\label{L4-2.1}
设 $\cC$ 是拟范畴。对形如
\[
\begin{tikzcd}[column sep=3.4em,row sep=2.4em]
& y \arrow[dr, "g"] & \\
x \arrow[ur, "f"] \arrow[rr, "h"'] & & z
\end{tikzcd}
\]
的 $2$-单形 $\sigma$，我们说 $\sigma$ \textbf{见证} $h$ 是 $f$ 与 $g$ 的复合。
\end{definition}

\begin{remark}\label{L4-2.2}
注意在上述定义中，态射 $h$ 不由 $f$ 与 $g$ 决定，$\sigma$ 也不由 $f$、$g$ 与 $h$ 决定。我们可以说：
\begin{center}
拟范畴中态射的复合不是具体的。
\end{center}
尽管如此，我们很快就会看到：只要正确理解“良定义”，由 $\cC$ 表示的 $\infty$-范畴中态射的复合就是良定义的。
\end{remark}

\section{态射之间的同伦}

\begin{definition}\label{L4-3.1}
设 $\cC$ 是拟范畴。对形如
\[
\begin{tikzcd}[column sep=3.4em,row sep=2.4em]
& y \arrow[dr, "\id"] & \\
x \arrow[ur, "f"] \arrow[rr, "g"'] & & y
\end{tikzcd}
\]
的 $2$-单形 $\sigma$，我们说 $\sigma$ 是从 $f$ 到 $g$ 的\textbf{左掐同伦}\footnote{此术语非标准。例如 [Rez22] 称之为右同伦。}，或\textbf{左掐 $2$-态射}。

对偶地，形如
\[
\begin{tikzcd}[column sep=3.4em,row sep=2.4em]
& x \arrow[dr, "f"] & \\
x \arrow[ur, "\id"] \arrow[rr, "g"'] & & y
\end{tikzcd}
\]
的 $2$-单形称为从 $f$ 到 $g$ 的\textbf{右掐同伦}，或\textbf{右掐 $2$-态射}。
\end{definition}

\yizhu{“掐”（pinched）取单形一端被捏合之意：左掐（右掐）$2$-单形的两条边中有一条是恒等态射，于是单形的两个顶点被掐合在一起。作者已注明此术语非标准，Rezk 称左掐同伦为右同伦，故阅读文献时须注意对应关系。\textbf{[译者按（无外部依据）]}}

\begin{exercise}\label{L4-3.2}
下列条件等价：
\begin{itemize}[leftmargin=2.2em,itemsep=0.1em]
\item 存在从 $f$ 到 $g$ 的左掐同伦；
\item 存在从 $f$ 到 $g$ 的右掐同伦。
\end{itemize}
此外，这些条件在从 $x$ 到 $y$ 的态射集合上定义了一个等价关系。提示：角嵌入 $\Lambda^3_1\to\Delta^3$ 与 $\Lambda^3_2\to\Delta^3$。
\end{exercise}

\begin{definition}\label{L4-3.3}
设 $\cC$ 是拟范畴，$f,g\colon x\to y$ 是态射。若它们满足上述练习中的条件，则称 $f$ \textbf{同伦于} $g$。
\end{definition}

\section{同伦范畴}

\begin{exercise}\label{L4-4.1}
设 $\cC$ 是拟范畴。证明下列构造定义了一个范畴 $h\cC$：
\begin{itemize}[leftmargin=2.2em,itemsep=0.1em]
\item $h\cC$ 的对象是 $\cC$ 的对象；
\item $h\cC$ 的态射是 $\cC$ 中态射的同伦类；
\item 在 $h\cC$ 中 $[f]$ 与 $[g]$ 的复合等于 $[h]$，当且仅当 $h$ 是 $\cC$ 中 $f$ 与 $g$ 的复合。
\end{itemize}
\end{exercise}

\begin{definition}\label{L4-4.2}
设 $\cC$ 是拟范畴。我们把上述练习中的范畴称为 $\cC$ 的\textbf{同伦范畴}，记作 $h\cC$。
\end{definition}

\begin{exercise}\label{L4-4.3}
我们有伴随对
\[
h\colon \mathrm{QCat} \rightleftarrows \Cat \colon N_{\bullet}.
\]
\end{exercise}

\begin{remark}\label{L4-4.4}
上述练习说明，$h\cC$ 是把由 $\cC$ 表示的 $\infty$-范畴中所有高阶态射强制为恒等态射后得到的范畴。注意这与放弃所有高阶态射并不相同。后一种构造是\textbf{邪恶的}，因为它违反同构原则。
\end{remark}

\nitem\label{L4-4.5} 顾名思义，拟范畴的同伦范畴与模型范畴的同伦范畴有关。我们将在以后的讲次中解释这种联系。

\section{同构}

\begin{definition}\label{L4-5.1}
设 $\cC$ 是拟范畴。若 $h\cC$ 中的态射 $[f]$ 是同构，则称态射 $f$ 是\textbf{同构}。
\end{definition}

\nitem\label{L4-5.2} 换言之，$f\colon x\to y$ 是同构，当且仅当存在态射 $g\colon y\to x$，使得 $g\circ f$ 同伦于 $\id_x$ 且 $f\circ g$ 同伦于 $\id_y$。

\begin{exercise}\label{L4-5.3}
拟范畴是 Kan 复形，当且仅当它的所有态射都是同构。
\end{exercise}

\begin{exercise}\label{L4-5.4}
对拟范畴 $\cC$，存在唯一的 Kan 复形 $\cC^{\simeq}\subset\cC$，称为 $\cC$ 的\textbf{核心}，它包含 $\cC$ 的全部同构。所得函子 $(-)^{\simeq}\colon\mathrm{QCat}\to\Kan$ 是嵌入函子 $\Kan\to\mathrm{QCat}$ 的右伴随。
\end{exercise}

\section{态射空间：想法}

\nitem\label{L4-6.1} 作为对 Joyal 提法的一次合理性检验，对拟范畴 $\cC$ 中的一对对象 $x,y$，我们应当构造一个 $\infty$-群胚 $\mathrm{Maps}_{\cC}(x,y)$，使得：
\begin{itemize}[leftmargin=2.2em,itemsep=0.1em]
\item $\mathrm{Maps}_{\cC}(x,y)$ 中的对象是从 $x$ 到 $y$ 的态射；
\item $\mathrm{Maps}_{\cC}(x,y)$ 中的态射是从 $x$ 到 $y$ 的态射之间的 $2$-态射；
\item $\cdots$
\end{itemize}
本讲给出三个代表这个 $\infty$-群胚的 Kan 复形构造（“代表”的含义见第 3 讲定义 8.3）。这三个 Kan 复形彼此弱同伦等价，但互不同构。

\nitem\label{L4-6.2} 在一个模型无关的、或公理化的 $(\infty,1)$-范畴理论\footnote{我并不知道这样的理论是否存在，因此下面的讨论仅起启发作用。}中，这三个构造应分别对应下述 $\infty$-群胚：
\begin{itemize}[leftmargin=2.6em,itemsep=0.15em,label={}]
\item[(B)] 纤维积
\[
\{x\}\times_{\Fun(\{0\},\cC)}\Fun([1],\cC)\times_{\Fun(\{1\},\cC)}\{y\},
\]
其中 $\Fun(\cD,\cC)$ 是从 $\cD$ 到 $\cC$ 的 $(\infty,1)$-范畴。
\item[(L)] 纤维积
\[
\cC_{x/}\times_{\cC}\{y\},
\]
其中 $\cC_{x/}$ 是 $\cC$ 在 $x$ 下的余切片 $(\infty,1)$-范畴。
\item[(R)] 纤维积
\[
\{x\}\times_{\cC}\cC_{/y},
\]
其中 $\cC_{/y}$ 是 $\cC$ 在 $y$ 上的切片 $(\infty,1)$-范畴。
\end{itemize}

\begin{exercise}\label{L4-6.3}
请说服自己：上述三个 $\infty$-群胚应当典范等价。
\end{exercise}

\section{函子拟范畴}

\nitem\label{L4-7.1} 为了用拟范畴实现 (B)，我们需要定义两个拟范畴之间的函子拟范畴。

\begin{definition}\label{L4-7.2}
称范畴 $\cC$ 是\textbf{笛卡尔闭}的，若它满足下列条件：
\begin{enumerate}[label=(\arabic*),leftmargin=2.4em,itemsep=0.1em]
\item $\cC$ 中有限积存在；
\item 对 $\cC$ 中任意对象 $X$ 与 $Y$，函子
\[
\Hom(X\times -,Y)\colon \cC^{\op}\to\Set
\]
余可表示，即存在对象 $Y^{X}$，称为\textbf{指数对象}，满足如下泛性质：
\[
\Hom(X\times Z,Y)\simeq\Hom(Z,Y^{X}).
\]
\end{enumerate}
\end{definition}

\begin{remark}\label{L4-7.3}
条件 (2) 等价于

(2$'$) 对任意对象 $X$，函子 $X\times -\colon \cC\to\cC$ 有右伴随。

事实上，这个右伴随由 $(-)^{X}\colon\cC\to\cC$ 给出。
\end{remark}

\begin{theoremdef}\label{L4-7.4}
范畴 $\Set_\Delta$ 是笛卡尔闭的。对单纯集 $X$ 与 $Y$，我们把指数对象记作
\[
\Fun(X,Y)\defeq Y^{X}\in\Set_\Delta.
\]
\end{theoremdef}

\begin{exercise}\label{L4-7.5}
$\Fun(X^{\op},Y^{\op})\simeq\Fun(X,Y)^{\op}$。
\end{exercise}

\begin{exercise}\label{L4-7.6}
证明脉函子 $N_{\bullet}\colon\Cat\to\Set_\Delta$ 与积和指数相容，即
\[
N_{\bullet}(\cC\times\cD)\xrightarrow{\ \simeq\ }N_{\bullet}(\cC)\times N_{\bullet}(\cD),
\qquad
N_{\bullet}(\Fun(\cC,\cD))\xrightarrow{\ \simeq\ }\Fun(N_{\bullet}(\cC),N_{\bullet}(\cD)).
\]
\end{exercise}

\begin{proposition}\label{L4-7.7}
（Ker.00TN）对 $X\in\Set_\Delta$ 与 $Y\in\Kan$，有 $\Fun(X,Y)\in\Kan$。
\end{proposition}

\begin{proposition}\label{L4-7.8}
（Ker.0066）对 $X\in\Set_\Delta$ 与 $Y\in\mathrm{QCat}$，有 $\Fun(X,Y)\in\mathrm{QCat}$。
\end{proposition}

\begin{corollary}\label{L4-7.9}
范畴 $\mathrm{QCat}$ 是笛卡尔闭的。
\end{corollary}

\begin{theoremdef}\label{L4-7.10}
（Ker.01JC）设 $\cC$ 是拟范畴，$x,y$ 是 $\cC$ 中的对象。则单纯集
\[
\Hom^{B}_{\cC}(x,y)\defeq\{x\}\times_{\Fun(\{0\},\cC)}\Fun(\Delta^1,\cC)\times_{\Fun(\{1\},\cC)}\{y\}
\]
是 Kan 复形，称为从 $x$ 到 $y$ 的\textbf{平衡态射复形}\footnote{此术语非标准。例如 Lurie 称之为态射空间。}。
\end{theoremdef}

\nitem\label{L4-7.11} 按定义，对检验单纯集 $K$，下述两个集合之间存在双射：
\begin{itemize}[leftmargin=2.2em,itemsep=0.1em]
\item 态射 $K\to\Hom^{B}_{\cC}(x,y)$；
\item 在 $K\times\{0\}$ 上取常值 $x$、在 $K\times\{1\}$ 上取常值 $y$ 的态射 $K\times\Delta^1\to\cC$。
\end{itemize}

\nitem\label{L4-7.12} 特别地，$\Hom^{B}_{\cC}(x,y)$ 中的一个 $1$-单形对应一个形如
\[
\begin{tikzcd}[column sep=3.2em,row sep=2.6em]
x \arrow[r, "f"] \arrow[d, "\id"'] \arrow[dr, dashed] & y \arrow[d, "\id"] \\
x \arrow[r, "g"'] & y
\end{tikzcd}
\]
的图 $\Delta^1\times\Delta^1\to\cC$。我们把这样的图称为从 $f$ 到 $g$ 的\textbf{平衡同伦}\footnote{此术语非标准。}，或\textbf{平衡 $2$-态射}。

\begin{exercise}\label{L4-7.13}
设 $\cC$ 是拟范畴，$f,g\colon x\to y$ 是态射。证明：$f$ 同伦于 $g$，当且仅当存在从 $f$ 到 $g$ 的平衡同伦。
\end{exercise}

\begin{exercise}\label{L4-7.14}
构造典范同构
\[
\Hom^{B}_{\cC}(x,y)^{\op}\xrightarrow{\ \simeq\ }\Hom^{B}_{\cC^{\op}}(y,x).
\]
\end{exercise}

\section{切片与余切片}

\nitem\label{L4-8.1} 为了用拟范畴实现 (L) 与 (R)，我们需要把切片/余切片范畴推广到拟范畴。

\begin{definition}\label{L4-8.2}
设 $\cC$ 与 $\cD$ 是范畴。$\cC$ 与 $\cD$ 的\textbf{联}（join）记作 $\cC\star\cD$，是如下定义的范畴：
\[
\Ob(\cC\star\cD)\defeq\Ob(\cC)\sqcup\Ob(\cD),
\]
\[
\Hom_{\cC\star\cD}(x,y)\defeq
\begin{cases}
\Hom_{\cC}(x,y), & x,y\in\cC,\\
\Hom_{\cD}(x,y), & x,y\in\cD,\\
\{\ast\}, & x\in\cC,\ y\in\cD,\\
\varnothing, & x\in\cD,\ y\in\cC.
\end{cases}
\]
\end{definition}

\begin{example}\label{L4-8.3}
$[i]\star[n-i-1]\simeq[n]$。
\end{example}

\nitem\label{L4-8.4} 上述定义可以推广到单纯集。

\begin{definition}\label{L4-8.5}
设 $X$ 与 $Y$ 是单纯集。$X$ 与 $Y$ 的\textbf{联}记作 $X\star Y$，是如下定义的单纯集：
\[
(X\star Y)([n])\defeq\bigsqcup_{-1\le i\le n}\bigl(X([i])\times Y([n-i-1])\bigr),
\]
其中我们采用约定 $[-1]=\varnothing$，从而 $X[-1]=Y[-1]=\{\ast\}$。在上述定义中，对 $[n]$ 的函子性由等价 $[i]\star[n-i-1]\simeq[n]$ 提供。
\end{definition}

\begin{example}\label{L4-8.6}
$\Delta^i\star\Delta^{n-i-1}\simeq\Delta^n$。
\end{example}

\begin{proposition}\label{L4-8.7}
（Ker.02QV）设 $X$ 与 $Y$ 是拟范畴。则联 $X\star Y$ 也是拟范畴。
\end{proposition}

\begin{definition}\label{L4-8.8}
设 $X$ 是单纯集。定义 $X$ 的\textbf{左锥}与\textbf{右锥}为
\[
X^{\triangleleft}\defeq\Delta^0\star X,
\qquad
X^{\triangleright}\defeq X\star\Delta^0.
\]
\end{definition}

\begin{definition}\label{L4-8.9}
设 $X$ 是单纯集，$x$ 是 $X$ 中的 $0$-单形。$X$ 在 $x$ 上的\textbf{切片单纯集}是如下定义的单纯集 $X_{/x}$：
\[
\Hom(Y,X_{/x})\simeq\Hom(Y^{\triangleright},X)\times_{\Hom(\Delta^0,X)}\{x\}.
\]
对偶地，$X$ 在 $x$ 下的\textbf{余切片单纯集}定义为满足如下泛性质的单纯集 $X_{x/}$：
\[
\Hom(Y,X_{x/})\simeq\Hom(Y^{\triangleleft},X)\times_{\Hom(\Delta^0,X)}\{x\}.
\]
\end{definition}

\begin{exercise}\label{L4-8.10}
证明：若 $X$ 是拟范畴，则 $X_{x/}$ 与 $X_{/x}$ 也是拟范畴。
\end{exercise}

\nitem\label{L4-8.11} 嵌入 $Y\to Y^{\triangleleft}$ 与 $Y\to Y^{\triangleright}$ 诱导遗忘态射
\[
X_{x/}\to X,
\qquad
X_{/x}\to X.
\]

\begin{theoremdef}\label{L4-8.12}
（Ker.01L0）设 $\cC$ 是拟范畴，$x,y$ 是 $\cC$ 中的对象。单纯集\footnote{下文中，$\{x\}$ 表示一个单纯集 $\Delta^0$，其 $0$-单形标为 $x$。}
\[
\Hom^{L}_{\cC}(x,y)\defeq\cC_{x/}\times_{\cC}\{y\},
\qquad
\Hom^{R}_{\cC}(x,y)\defeq\{x\}\times_{\cC}\cC_{/y}
\]
是 Kan 复形，称为从 $x$ 到 $y$ 的\textbf{左/右掐态射复形}。
\end{theoremdef}

\nitem\label{L4-8.13} 按定义，对检验单纯集 $K$，下述两个集合之间存在双射：
\begin{itemize}[leftmargin=2.2em,itemsep=0.1em]
\item 态射 $K\to\Hom^{L}_{\cC}(x,y)$；
\item 在顶点 $\Delta^0\subset K^{\triangleleft}$ 上取 $x$、在底 $K\subset K^{\triangleleft}$ 上取常值 $y$ 的态射 $K^{\triangleleft}\to\cC$。
\end{itemize}
对偶地，下述两个集合之间也存在双射：
\begin{itemize}[leftmargin=2.2em,itemsep=0.1em]
\item 态射 $K\to\Hom^{R}_{\cC}(x,y)$；
\item 在顶点 $\Delta^0\subset K^{\triangleright}$ 上取 $y$、在底 $K\subset K^{\triangleright}$ 上取常值 $x$ 的态射 $K^{\triangleright}\to\cC$。
\end{itemize}
特别地，$\Hom^{L}_{\cC}(x,y)$ 与 $\Hom^{R}_{\cC}(x,y)$ 中的态射分别代表 $\cC$ 中的左/右掐 $2$-态射。

\begin{exercise}\label{L4-8.14}
构造典范同构
\[
\Hom^{L}_{\cC}(x,y)^{\op}\xrightarrow{\ \simeq\ }\Hom^{R}_{\cC^{\op}}(y,x).
\]
\end{exercise}

\section{态射空间}

\nitem\label{L4-9.1} 对检验单纯集 $K$，有自然映射
\[
K^{\triangleleft}\leftarrow K\times\Delta^1\to K^{\triangleright},
\]
其中向左（相应地，向右）的态射把 $K\times\{0\}$（相应地，$K\times\{1\}$）坍缩到左（相应地，右）锥的顶点。由 \S\ref{L4-8.13} 与 \S\ref{L4-7.11}，我们得到态射
\begin{equation}\label{L4-9.1eq}
\Hom^{L}_{\cC}(x,y)\to\Hom^{B}_{\cC}(x,y)\leftarrow\Hom^{R}_{\cC}(x,y).
\end{equation}

\begin{proposition}\label{L4-9.2}
（Ker.01L5）\eqref{L4-9.1eq} 中的态射是弱同伦等价。
\end{proposition}

\nitem\label{L4-9.3} 由此可见，由 \eqref{L4-9.1eq} 中 Kan 复形所代表的 $\infty$-群胚/空间典范等价。

\begin{definition}\label{L4-9.4}
设 $\cC$ 是拟范畴，$x,y$ 是 $\cC$ 中的对象。设
\[
\mathrm{Maps}_{\cC}(x,y)\in\hGrpd_\infty\defeq\Set_\Delta[W^{-1}]
\]
是由 \eqref{L4-9.1eq} 代表的对象\footnote{其确切含义见 \S11。}。我们称它为从 $x$ 到 $y$ 的\textbf{映射空间}。
\end{definition}

\yizhu{这里的 $W$ 指 $\Set_\Delta$ 的 Quillen 模型结构中的弱等价类，$\Set_\Delta[W^{-1}]$ 是把 $W$ 形式可逆化所得的同伦范畴，即 $\infty$-群胚的范畴（见第 3 讲）。\textbf{[来源：B]}}

\nitem\label{L4-9.5} 不管下面这句话是什么意思，我们把 $\infty$-群胚 $\mathrm{Maps}_{\cC}(x,y)$ 中的态射称为由 $\cC$ 表示的 $\infty$-范畴中的 \textbf{$2$-态射}。

\begin{exercise}\label{L4-9.6}
构造典范同构
\[
\mathrm{Maps}_{\cC}(x,y)\xrightarrow{\ \simeq\ }\mathrm{Maps}_{\cC^{\op}}(x,y).
\]
提示：第 3 讲练习 6.5。
\end{exercise}

\section{$\infty$-范畴中的复合}

\nitem\label{L4-10.1} 下述结果属于 Joyal，它是组合性的，但很基本。证明见 Ker.0079。

\begin{theorem}\label{L4-10.2}
单纯集 $X$ 是拟范畴，当且仅当
\[
\Fun(\Delta^2,X)\to\Fun(\Lambda^2_1,X)
\]
是 $\Set_\Delta$ 中的无环 Kan 纤维化。
\end{theorem}

\nitem\label{L4-10.3} 设 $\cC$ 是拟范畴。按定义，$\Fun(\Lambda^2_1,\cC)$ 中的一个 $0$-单形是 $\cC$ 中的一对可复合态射：
\[
\begin{tikzcd}[column sep=3.4em,row sep=2.4em]
& y \arrow[dr, "g"] & \\
x \arrow[ur, "f"] \arrow[rr] & & z
\end{tikzcd}
\]
而把这个 $0$-单形提升到 $\Fun(\Delta^2,\cC)$ 的一个提升，就是一个 $2$-单形 $\sigma$
\begin{equation}\label{L4-10.1eq}
\begin{tikzcd}[column sep=3.4em,row sep=2.4em]
& y \arrow[dr, "g"] & \\
x \arrow[ur, "f"] \arrow[rr, "h"'] & & z
\end{tikzcd}
\end{equation}
见证 $h$ 是 $f$ 与 $g$ 的复合。现在，上述定理说明这种提升构成的单纯集是弱可缩的 Kan 复形。换言之，$f$ 与 $g$ 的复合——当理解为一个 $2$-单形时——在可缩的选择空间意义下良定义。由下面的口号~\ref{L4-11.4}，我们可以说：
\begin{center}
\emph{复合 $g\circ f$ 作为映射 $\infty$-群胚 $\mathrm{Maps}_{\cC}(x,y)$ 中的对象是良定义的。}
\end{center}

\yizhu{原文 \S10.3 与 \S11.1 两处均把复合 $g\circ f$（$x\to y\to z$）说成 $\mathrm{Maps}_{\cC}(x,y)$ 中的对象，\S10.7 第 (1) 小问又写作 $\Fun(\Lambda^2_1,X)$；按上下文应分别为 $\mathrm{Maps}_{\cC}(x,z)$ 与 $\Fun(\Lambda^2_1,\cC)$。译文均照录原文。\textbf{[译者按（无外部依据）]}}

\nitem\label{L4-10.4} 为得到上述结果，关键是把见证 $2$-单形作为复合的数据的一部分纳入进来。我们可以说：

\begin{slogan}\label{L4-10.5}\centering\par
在 $\infty$-范畴中，复合是结构而非性质。
\end{slogan}

\nitem\label{L4-10.6} 由于这一点相当重要，我们换一种说法重述。在 $\infty$-范畴中，当我们说 $h$ 是 $f$ 与 $g$ 的复合时，我们总是指我们给出了一个交换图 \eqref{L4-10.1eq}。否则，“复合”一词没有意义。

\begin{exercise}\label{L4-10.7}
另一种为上述断言提供依据的方式如下。设 $\cC$ 是拟范畴。
\begin{enumerate}[label=(\arabic*),leftmargin=2.4em,itemsep=0.1em]
\item 用定理~\ref{L4-10.2} 推出
\[
\Fun(\Delta^2,\cC)\times_{\cC^{\times 3}}\{(x,y,z)\}\to\Fun(\Lambda^2_1,X)\times_{\cC^{\times 3}}\{(x,y,z)\}
\]
是 $\Set_\Delta$ 中的无环 Kan 纤维化。
\item 把上述态射的目标等同于
\[
\Hom^{B}_{\cC}(x,y)\times\Hom^{B}_{\cC}(y,z).
\]
\item 在 $\Set_\Delta$ 中构造态射
\[
\Fun(\Delta^2,\cC)\times_{\cC^{\times 3}}\{(x,y,z)\}\to\Hom^{B}_{\cC}(x,z).
\]
\item 在 $\hGrpd_\infty\defeq\Set_\Delta[W^{-1}]$ 中构造态射
\[
\Hom^{B}_{\cC}(x,y)\times\Hom^{B}_{\cC}(y,z)\to\Hom^{B}_{\cC}(x,z).
\]
\item 证明 (4) 中的复合满足严格的单位律与结合律公理。
\end{enumerate}
\end{exercise}

\begin{exercise}\label{L4-10.8}
在 $\infty$-范畴 $\cC$ 中，与同构复合会诱导映射空间之间的等价。换言之，对拟范畴 $\cC$：
\begin{enumerate}[label=(\arabic*),leftmargin=2.4em,itemsep=0.1em]
\item 若 $f\colon x\to y$ 是同构，则它诱导 Kan 复形之间的弱同伦等价
\[
-\circ f\colon\Hom^{B}_{\cC}(y,z)\xrightarrow{\ \sim\ }\Hom^{B}_{\cC}(x,z).
\]
\item 若 $g\colon y\to z$ 是同构，则它诱导 Kan 复形之间的弱同伦等价
\[
g\circ -\colon\Hom^{B}_{\cC}(x,y)\xrightarrow{\ \sim\ }\Hom^{B}_{\cC}(x,z).
\]
\end{enumerate}
\end{exercise}

\section{对象何时良定义？}

\nitem\label{L4-11.1} 本节旨在回答下述问题：
\begin{center}
\emph{为什么 $g\circ f$ 作为 $\infty$-群胚 $\mathrm{Maps}_{\cC}(x,y)$ 中的对象是良定义的？}
\end{center}

\nitem\label{L4-11.2} 首先说明：我们并不断言 $g\circ f$ 是拟范畴 $\cC$ 中良定义的 $1$-单形，也不断言它是 \eqref{L4-9.1eq} 中任何一个 Kan 复形中良定义的 $0$-单形。我们断言的是：它作为由 $\cC$ 建模的 $\infty$-范畴中的态射是良定义的，等价地，作为由 \eqref{L4-9.1eq} 建模的 $\infty$-群胚中的对象是良定义的。既然我们已在练习~\ref{L4-10.7} 中勾画了这一断言的数学证明，下面我们给出一个更具形而上学/观念形态色彩的论证。

\nitem\label{L4-11.3} 在经典范畴论中，\textbf{构造}常常意味着
\begin{itemize}[leftmargin=2.2em,itemsep=0.1em]
\item 指定给定集合中的一个元素。
\end{itemize}
然而在 $\infty$-范畴论中，我们希望追踪对象之间的自同构，以及自同构之间的自同构，等等。因此构造应当意味着
\begin{itemize}[leftmargin=2.2em,itemsep=0.1em]
\item 指定给定 $\infty$-群胚中的一个对象，
\end{itemize}
或等价地，
\begin{itemize}[leftmargin=2.2em,itemsep=0.1em]
\item 指定一个同伦型中的一个点。
\end{itemize}
注意我们无法区分一个点的同伦型与任何可缩空间的同伦型。因此我们可以说：

\begin{slogan}\label{L4-11.4}\centering\par
在 $\infty$-群胚中，一个对象是良定义的，如果它在可缩的选择空间意义下良定义\footnote{这一哲学的极端推论是所谓的数学的单值基础。}。
\end{slogan}

\nitem\label{L4-11.5} 从现在起，
\begin{center}
本质唯一 $=$ 在可缩的选择空间意义下唯一。
\end{center}

\appendixsection{逆态射}

\nitem\label{L4-A.1} \S\ref{L4-1.3} 中的提升性质实际上蕴含 $f$ 是同构（而不只是有左逆），只要单纯集是拟范畴。更确切地说：

\begin{exercise}\label{L4-A.2}
设 $\cC$ 是拟范畴，$f$ 是态射。假设图
\[
\begin{tikzcd}[column sep=3.4em,row sep=2.6em]
\Lambda^n_0 \arrow[r, "\phi"] \arrow[d] & X \\
\Delta^n \arrow[ur, dashed] &
\end{tikzcd}
\qquad n\ge 2
\]
对任何满足 $\phi|_{\Delta^{\{0,1\}}}=f$ 的 $\phi$ 都有解。证明 $f$ 是同构。
\end{exercise}

\nitem\label{L4-A.3} 下述练习说明：同构 $f$ 的左逆、右逆与逆都是良定义的，即在可缩的选择空间意义下定义。不过，须小心这些词的含义。

\begin{exercise}\label{L4-A.4}
设 $\cC$ 是拟范畴，$f$ 是态射。找出下列说法的正确含义：
\begin{itemize}[leftmargin=2.2em,itemsep=0.1em]
\item $2$-单形 $\sigma$ 见证 $g$ 是 $f$ 的左/右逆。
\end{itemize}
若 $f$ 是同构，仿照 \S10 证明 $f$ 的左/右逆是良定义的。
\end{exercise}

\begin{exercise}\label{L4-A.5}
设 $\cC$ 是拟范畴，$f$ 是同构。求一个单纯集 $I$，使得下述说法
\begin{itemize}[leftmargin=2.2em,itemsep=0.1em]
\item 态射 $I\to\cC$ 见证 $g$ 是 $f$ 的逆
\end{itemize}
是正确的概念。仿照 \S10 证明 $f$ 的逆是良定义的。提示：$I$ 应当是某个普通范畴的脉。
\end{exercise}

\begin{exercise}\label{L4-A.6}
请为下述说法赋予意义：在 $\infty$-范畴中，同构的左逆与右逆典范等价。
\end{exercise}

\subsection*{参考文献}
\par\noindent\hangindent=5.2em\hangafter=1 [Rez22]\hspace{0.6em}Charles Rezk. \emph{Introduction to quasicategories}. Lecture Notes for course at University of Illinois at Urbana-Champaign, 2022.
